\documentclass[11pt,a4paper]{article}

\usepackage[utf8]{inputenc}
\usepackage[T1]{fontenc}
\usepackage[ngerman]{babel}
\usepackage{amsmath,amssymb,amsthm}
\usepackage{array,booktabs}
\usepackage{geometry}
\usepackage{enumitem}

\geometry{margin=2.5cm}

\newcommand{\tr}{\operatorname{tr}}
\newcommand{\rg}{\operatorname{rg}}
\newcommand{\T}{^{\!\top}}
\newcommand{\adj}{^{\ast}}

\newcolumntype{L}{>{\raggedright\arraybackslash}p{0.46\linewidth}}

\title{Die Invariante~19:\\[0.3em]
  Sieben unzertrennliche Dualpaare\\
  und der lineare Fixpunkt $P_{19}$}
\author{}
\date{}

\begin{document}
\maketitle

\section*{Aufbau}

Das System besteht aus \textbf{vierzehn Axiomen} $A_1,\dots,A_{14}$,
die zu \textbf{sieben unzertrennlichen Dualpaaren} geordnet sind:
\[
  (A_1,A_{14}),\;(A_2,A_{12}),\;(A_3,A_8),\;(A_4,A_{10}),\;
  (A_5,A_7),\;(A_6,A_{11}),\;(A_9,A_{13}).
\]
Jedes Paar besitzt zwei Seiten:

\begin{description}[leftmargin=2em,style=nextline]
  \item[Wirkungsaspekt (Operator).] Was \emph{tut} der Operator? Eine
    Operation, Transformation oder Wirkung auf einem Raum.
  \item[Strukturaspekt (Raum).] Welche \emph{Eigenschaft} ist dem Raum
    bzw.\ dem Operator selbst eigen, unabhängig von einer einzelnen
    Wirkung? Das Reaktiv-, Reflexions- oder Gleichgewichtsmoment.
\end{description}

Der \textbf{Lineare Fixpunkt} $P_{19}$ ist \emph{nicht} ein achtes Paar,
sondern der simultane Konvergenzpunkt aller sieben Paare: jenes Objekt,
auf dem jede der vierzehn Axiom-Operationen den Eigenwert $1$ realisiert.
Wir notieren dies einheitlich
\[
  \boxed{\;A_k(P_{19}) = 1\cdot P_{19}\quad\text{für alle}\quad
    k\in\{1,\dots,14\}\;}
\]
Auf $P_{19}$ kollabiert in jedem Paar die Unterscheidung zwischen
Wirkungs- und Strukturaspekt --- Operator und Raum sind dort
ununterscheidbar.

\paragraph{Notationskonvention.}
Damit das Symbol $P$ \emph{ausschließlich} für den Fixpunkt $P_{19}$
reserviert bleibt, verwenden wir:
\begin{itemize}[leftmargin=2em]
  \item $S\in\mathrm{GL}(n)$ für Ähnlichkeitstransformationen,
  \item $\Pi_i$ für Spektralprojektoren,
  \item $Q\in O(n)$ für orthogonale Operatoren,
  \item $A\adj$ für den adjungierten Operator (reell: $A\T$;
        komplex: $A^{\dagger}$).
\end{itemize}

\bigskip
\section*{Die sieben Dualpaare}

\subsection*{I.\ Grenze vs.\ Vollständigkeit}
\begin{tabular}{@{}LL@{}}
  \toprule
  \textbf{$A_1$ -- Singularität (Kern)} &
  \textbf{$A_{14}$ -- Vollständigkeit (Spektralzerlegung)} \\
  \midrule
  $\exists\, v\neq 0:\; A v = 0$, d.h.\ $\ker(A)\neq\{0\}$. &
  $\displaystyle \sum_i \Pi_i = I$ mit Spektralprojektoren $\Pi_i$. \\
  Das Verschwinden im Nullraum. &
  Die identische Auflösung in alle Eigenrichtungen. \\
  \bottomrule
\end{tabular}

\subsection*{II.\ Aktion vs.\ Reflexion}
\begin{tabular}{@{}LL@{}}
  \toprule
  \textbf{$A_2$ -- Kopplung (Tensorprodukt)} &
  \textbf{$A_{12}$ -- Rückkopplung (Resolvente)} \\
  \midrule
  $V\otimes W$ als Interaktionsraum zweier Systeme. &
  $\mathcal{L}\{e^{At} x_0\}(s) = (sI-A)^{-1}\, x_0$. \\
  Die offene Kopplung nach außen. &
  Der geschlossene Kreis der Selbstwirkung. \\
  \bottomrule
\end{tabular}

\subsection*{III.\ Ähnlichkeit vs.\ Selbstadjungiertheit}
\begin{tabular}{@{}LL@{}}
  \toprule
  \textbf{$A_3$ -- Spur-Invarianz} &
  \textbf{$A_8$ -- Selbstadjungiertheit} \\
  \midrule
  $\tr(S^{-1} A S) = \tr(A)$ für $S\in\mathrm{GL}(n)$. &
  $A\adj = A$ \quad
    (reell: $A\T = A$, komplex: $A^{\dagger} = A$). \\
  Formkonstanz unter Basiswechsel. &
  Symmetrie als innere Eigenschaft. \\
  \bottomrule
\end{tabular}

\subsection*{IV.\ Takt vs.\ Resonanz}
\begin{tabular}{@{}LL@{}}
  \toprule
  \textbf{$A_4$ -- Frequenz (Eigenwert)} &
  \textbf{$A_{10}$ -- Resonanz (Entartung)} \\
  \midrule
  $A v = \lambda v$, $\ v\neq 0$. &
  $\lambda_i = \lambda_j\;\Longrightarrow\;
   \dim\ker(A-\lambda I) > 1$. \\
  Diskreter Eigentakt einer Richtung. &
  Kopplung mehrerer Richtungen im selben Takt. \\
  \bottomrule
\end{tabular}

\subsection*{V.\ Wirbel vs.\ Quelle}
\begin{tabular}{@{}LL@{}}
  \toprule
  \textbf{$A_5$ -- Rotation (Orthogonalität)} &
  \textbf{$A_7$ -- Feldquelle (Divergenz/Spur)} \\
  \midrule
  $Q\in O(n):\; \lVert Q v\rVert = \lVert v\rVert$. &
  $\tr(A) = \sum_i a_{ii}$ als Quellstärke (Divergenz). \\
  Drehung erhält Länge, erzeugt aber keinen Fluss. &
  Stationäre Quell- oder Senkenstärke. \\
  \bottomrule
\end{tabular}

\subsection*{VI.\ Gefälle vs.\ Akkumulation}
\begin{tabular}{@{}LL@{}}
  \toprule
  \textbf{$A_6$ -- Potential (quadratische Form)} &
  \textbf{$A_{11}$ -- Integration (Skalarprodukt)} \\
  \midrule
  $\Phi(x) = x\T A\, x$. &
  $\langle u, v\rangle = \displaystyle\int u\cdot v\,\mathrm{d}\mu$. \\
  Lokales Gefälle einer Energiefläche. &
  Globale Aufsummierung über das Maß. \\
  \bottomrule
\end{tabular}

\subsection*{VII.\ Stufen vs.\ Kontinuum}
\begin{tabular}{@{}LL@{}}
  \toprule
  \textbf{$A_9$ -- Quantisierung (Dimension)} &
  \textbf{$A_{13}$ -- Dichte (Rang-Metrik)} \\
  \midrule
  $\dim(V) = n \in \mathbb{N}$. &
  $\displaystyle \rho \;=\; \frac{\rg(A)}{\min(n,m)} \in [0,1]$
   für $A\in\mathbb{R}^{n\times m}$. \\
  Diskrete Stufung des Raums. &
  Stetiges Maß der inneren Belegung. \\
  \bottomrule
\end{tabular}

\bigskip
\section*{Der lineare Fixpunkt $P_{19}$}

\begin{quote}
\itshape
$P_{19}$ ist der einzige Punkt, an dem alle vierzehn Axiome simultan zum
Eigenwert~$1$ konvergieren --- d.h.\ der Schnitt der vierzehn Eins-Eigenräume
\[
  P_{19} \;\in\; \bigcap_{k=1}^{14}\ \{\,x : A_k(x) = 1\cdot x\,\}.
\]
\end{quote}

Eigenschaften:
\begin{enumerate}[leftmargin=2em]
  \item \textbf{Simultanität.} Innerhalb jedes der sieben Dualpaare wirken
    Wirkungs- und Strukturaspekt auf $P_{19}$ identisch. Damit ist die
    Unterscheidung Operator/Raum aufgehoben.
  \item \textbf{Linearität.} Aus $A(P_{19}) = 1\cdot P_{19}$ folgt
    $A^n(P_{19}) = P_{19}$ für alle $n\in\mathbb{N}$: $P_{19}$ ist unter
    Iteration stabil.
  \item \textbf{Eindeutigkeit als Schnittpunkt.} $P_{19}$ ist nicht ein
    achtes Paar, sondern derjenige Punkt, in dem die sieben Paare zusammen-
    fallen --- die Invariante~19.
\end{enumerate}

\end{document}
